% §18.2 Bloom-Filter-Based Constructions
%
% Monograph register: the per-document and inverted-index constructions,
% each as a cipher map, FNR-zero structural property, space/query cost,
% worked example. Tightened.
% FIX: the conjunctive-query FPR reasoning. The draft said a conjunctive
% false positive "requires both individual filters to falsely match"
% (implying eps^2). The headline bound "<= eps per document" is correct,
% but the reason is wrong: a document containing w1 but not w2 (which is
% NOT in the conjunction) is a false positive when only the w2-test
% falsely matches, at rate eps. The rate is the ch6 three-region sum
% (eq:intersection-fpr): eps in the missing-one-keyword regions, ~eps^2
% only for documents missing both. Corrected.
% Terminology: "match obliviously" / "oblivious to the underlying
% keywords" kept (the oblivious-to-X idiom).
% Cross-refs thm:composition, def:bloom-filter, sec:bloom-from-scratch,
% thm:bloom-optimal-k, sec:union-of-sets, prop:intersection-fpr-fnr,
% eq:intersection-fpr, def:cipher-map, def:encrypted-search. DEFINED:
% sec:bloom-based, def:per-document-bloom, def:inverted-index-bloom.

\section{Bloom-Filter-Based Constructions}
\label{sec:bloom-based}

The cipher-map framing of \Cref{sec:encrypted-search-problem} identified
encrypted search with a family of approximate membership tests on
encrypted documents. A construction needs an approximate-set structure
recording each document's keywords compactly, an encryption scheme
letting the server test a query token against the stored representation
without learning the plaintext keyword, and a query-token format the
client computes privately and the server matches obliviously. The Bloom
filter of \Cref{def:bloom-filter} solves the first; a standard symmetric
scheme (a one-time pad keyed by a pseudorandom function) the second; a
deterministic keyed function of the keyword the third. The two
construction families differ in what counts as a ``document'' in the
Bloom sense: the first puts a filter inside each document, the second
inside each keyword.

\subsection*{The two constructions}

\begin{definition}[Per-document Bloom-filter encrypted search]
\label{def:per-document-bloom}
For each document $d_i$ the client computes a Bloom filter $B_i$ over its
keywords, encrypts it as $\mathsf{Enc}(B_i)$ under a key the server does
not hold, and uploads $\{(i, \mathsf{Enc}(B_i))\}$. To search for $w$ the
client sends a query token $\tau_w$ (a deterministic function of $w$ and
its secret key); the server tests $\tau_w$ against each
$\mathsf{Enc}(B_i)$ and returns the identifiers for which the test
succeeds.
\end{definition}

\begin{definition}[Inverted-index Bloom-filter encrypted search]
\label{def:inverted-index-bloom}
For each keyword $w$ the client computes a Bloom filter $B_w$ over the
identifiers $\{i : w \in d_i\}$, encrypts it, and uploads
$\{(\tau_w, \mathsf{Enc}(B_w))\}$ indexed by token. To search for $w$ the
client sends $\tau_w$; the server returns the matching encrypted filter,
which the client decrypts to recover an approximate identifier set.
\end{definition}

The per-document construction asks, for each document, ``does this query
match?'', and has $|D|$ answers per query; the inverted-index
construction asks ``which documents match?'', with $|W|$ index entries
and a shorter query path (one lookup, one encrypted blob, the rest done
client-side). The encryption layer does two jobs in both: it hides the
bit pattern from the server, and it makes the match oblivious to the
underlying keywords, since $\tau_w$ carries $w$ only through a keyed
function the server cannot invert.

\subsection*{Each is a cipher map, and FNR-zero}

In the per-document construction the plaintext function is $f_i(w) =
\mathbf{1}[w \in d_i]$: the token replaces $w$, the encrypted filter
replaces $d_i$, and the server-side match approximates $f_i$ up to the
filter's $\fprate$, always reporting a match when $w \in d_i$ ($\fnrate =
0$) and reporting one with probability at most $\fprate$ when $w \notin
d_i$. The error profile is $(\fprate, 0)$, satisfying
\Cref{def:cipher-map}. The inverted-index construction realizes $f(w) =
\{i : w \in d_i\}$ with the same profile. Both have $\fnrate = 0$ by
structural inheritance from the Bloom filter, whose union analysis
(\Cref{sec:union-of-sets}) showed that ORing bits can only add matches,
never remove them, so the per-document filter never misses an inserted
keyword.

The chapter 3 composition theorem (\Cref{thm:composition}) extends this
to multi-keyword conjunctions, with the three-region care chapter 6
required for intersection. A client querying for documents containing
both $w_1$ and $w_2$ intersects the two returned identifier sets. The
result has $\fnrate = 0$ (a true match, containing both keywords, is
returned by both queries and survives). Its false-positive rate per
document is bounded by $\fprate$, but not for the reason a bare product
would suggest. A document outside the conjunction is a false positive
when the test for a keyword it lacks falsely matches: a document
containing $w_1$ but not $w_2$ is returned whenever the $w_2$-test errs,
at rate $\fprate$ (the $w_1$-test is correct, not a false match), and only
a document containing neither keyword needs both tests to err, at the
lower rate $\approx \fprate^2$. By the three-region intersection rate
(\Cref{prop:intersection-fpr-fnr}, \eqref{eq:intersection-fpr}) the
conjunction's FPR is the partition-weighted sum, bounded by $\fprate$ and
met by documents missing exactly one of the two keywords. Deep
conjunctions are highly precise on a corpus of true non-matches and only
modestly precise when documents commonly contain one queried keyword but
not the other.

A Bloom-based encrypted search therefore never \emph{loses} a real match:
every document truly containing the keyword appears. What varies is how
many \emph{additional} documents appear as false positives, the
operational price the client pays, exactly the price chapter 5 set for
the underlying filter.

\subsection*{Space and query cost}

For the per-document construction with $|D|$ documents and $W_d$ keywords
each, the chapter 5 optimal-$k$ formula (\Cref{thm:bloom-optimal-k}) gives
$m_d / W_d = (\log_2 e)\log_2(1/\fprate) \approx 1.44 \log_2(1/\fprate)$
bits per item, so total storage is
\[
  S_{\text{per-doc}} \;\approx\; 1.44 \cdot |D| \cdot W_d \cdot
  \log_2(1/\fprate),
\]
at $|D|$ server-side matches per query. The inverted-index construction
with vocabulary $|W|$ and average per-keyword document count $|D_w|$ costs
$S_{\text{inv}} \approx 1.44 \cdot |W| \cdot |D_w| \cdot \log_2(1/\fprate)$
at one server lookup per query. Since $\sum_d W_d = \sum_w |D_w|$ is the
total keyword-document incidence, the two cost the same in storage to
leading order; they differ in query cost by a factor of $|D|$ at the
server (the per-document construction tests every filter, the
inverted-index does one lookup), at the price of pushing the membership
test to the client. Incremental update favors the per-document layout
(adding a document is local) while query latency favors the inverted
index, so production deployments routinely keep both, the per-document
filters as the primary store and the inverted index as a query
accelerator rebuilt lazily.

\subsection*{A worked example}

Take $|D| = 10{,}000$ documents, $W_d = 50$ keywords each, $\fprate =
0.01$. The per-document filter has $1.44 \log_2(100) \approx 9.59$ bits
per item, so $m_d \approx 50 \cdot 9.59 \approx 480$ bits, total
$10{,}000 \cdot 480 = 4.8 \times 10^6$ bits $= 0.6$ MB. Exact word-set
storage at 50 keywords of 8 bytes each is $10{,}000 \cdot 50 \cdot 8 = 4$
MB, about seven times more, and the Bloom representation also encrypts a
fixed-size object regardless of keyword length. The inverted-index
version needs the same $\sum_d W_d = 5 \times 10^5$ incidence pairs at
$9.59$ bits each, the same $0.6$ MB: the construction style changes the
layout and the query path, not the storage cost. At $\fprate = 0.01$ over
this corpus the expected spurious matches per query are $0.01 \cdot |D| =
100$, against perhaps $|D| \cdot \pi = 10$ true matches at prevalence
$\pi = 0.001$; the response is dominated by false positives, which the
client filters out by some means external to the encrypted-search layer.
The next section reads those false positives as the source of the
system's confidentiality rather than as a defect.
